% 双产出收益最大化案例
% 案例统一前言（与前三讲书稿共用版式）
\newcommand{\lectureno}{案例}
\newcommand{\lectureshort}{双产出收益最大化}
\newcommand{\booktitle}{微观经济学理论：生产者选择与对偶}
% ============================================================
%  微观经济学理论课程讲义 —— 统一课程前言（四讲共用）
%  版式参考 Nolan Miller, Notes on Microeconomic Theory (SWP/notes3 风格)
%  编译：LuaLaTeX（两遍）
%  各讲 preamble.tex 只需设定 \lectureno / \lectureshort 后 \input 本文件。
% ============================================================
\documentclass[11pt,a4paper,fontset=none,oneside,openany]{ctexbook}

% ---------- 版式 ----------
\usepackage{geometry}
\geometry{
  a4paper,
  top=2.8cm, bottom=2.8cm,
  left=2.6cm, right=2.6cm,
  headsep=0.85cm, footskip=1.25cm
}
\setlength{\headheight}{14pt}

% ---------- 字体（全部矢量嵌入，保证 PDF 高清晰度）----------
% 注意：本机 macOS 已移除 Kaiti SC，CJK 斜体统一回退到 Songti SC；
%        Songti SC / Heiti SC 经 OSFONTDIR 指向 Supplemental 后由 luaotfload 解析。
\usepackage{fontspec}
\setmainfont{TeX Gyre Termes}[Ligatures=TeX]
\setsansfont{TeX Gyre Heros}[Ligatures=TeX,Scale=MatchLowercase]
\setmonofont{TeX Gyre Cursor}[Scale=MatchLowercase]
\setCJKmainfont{Songti SC}[BoldFont={Songti SC Bold},ItalicFont={Songti SC}]
\setCJKsansfont{Heiti SC}[BoldFont={Heiti SC}]

% ---------- 数学 ----------
\usepackage{amsmath}
\usepackage{mathtools}
\usepackage{amsthm}
\usepackage{unicode-math}
\setmathfont{TeX Gyre Termes Math}
\allowdisplaybreaks[2]
\setlength{\abovedisplayskip}{9pt plus 3pt minus 3pt}
\setlength{\belowdisplayskip}{9pt plus 3pt minus 3pt}

% ---------- 定理类环境（书稿格式：按章编号，如“定理 1.1”）----------
\theoremstyle{plain}
\newtheorem{theorem}{定理}[chapter]
\newtheorem{proposition}[theorem]{命题}
\newtheorem{lemma}[theorem]{引理}
\newtheorem{corollary}[theorem]{推论}
\theoremstyle{definition}
\newtheorem{definition}[theorem]{定义}
\newtheorem{condition}[theorem]{条件}
\newtheorem{example}[theorem]{例}
\theoremstyle{remark}
\newtheorem{remark}[theorem]{注}
% 公式、图、表按章编号（书稿格式：式 (1.1)、图 1.1、表 1.1）
\numberwithin{equation}{chapter}
\numberwithin{figure}{chapter}
\numberwithin{table}{chapter}

% 定理类环境标题后的标点：amsthm 默认加英文 "."，此处统一改为中文句号。
% 原理：amsthm 的 \@thm 每次开环境都执行 \thm@headpunct{.}，而样式宏在其之后展开，
% 故把 \thm@headpunct{。} 追加到三个样式宏末尾即可覆盖（用 \AtBeginDocument 以免被 ctex 覆盖）。
\makeatletter
\AtBeginDocument{%
  \let\WBth@plain\th@plain
  \let\WBth@definition\th@definition
  \let\WBth@remark\th@remark
  \def\th@plain{\WBth@plain\thm@headpunct{。}}%
  \def\th@definition{\WBth@definition\thm@headpunct{。}}%
  \def\th@remark{\WBth@remark\thm@headpunct{。}}%
}
\makeatother

% 无编号小标题（仿 notes3 的 “A Note on ...”）
\newcommand{\notehead}[1]{\par\medskip\noindent{\bfseries #1}\par\smallskip}

% ---------- 表格 ----------
\usepackage{booktabs}
\usepackage{array}
\renewcommand{\arraystretch}{1.25}
\setlength{\tabcolsep}{8pt}

% ---------- 图形 ----------
\usepackage{graphicx}
\usepackage{tikz}
\usepackage{pgfplots}
\pgfplotsset{compat=1.18}
\usetikzlibrary{arrows.meta,positioning,shapes.geometric,calc,patterns,decorations.pathreplacing,decorations.markings}
\usepackage{float}
\usepackage{caption}
\captionsetup{font=small,labelfont=bf,labelsep=quad,skip=4pt}

% 统一图形配色（黑白模式：黑/灰梯度，供所有 TikZ 图复用）
% —— 全局黑白化：原彩色调色板统一改为黑/灰，保证打印与投影均可读 ——
\definecolor{AxisBlue}{RGB}{0,0,0}        % 轴线 / 主曲线：黑
\definecolor{SoftBlue}{RGB}{130,130,130}  % 节点填充：中灰
\definecolor{LightRed}{RGB}{224,224,224}  % 浅填充：浅灰
\definecolor{DeepBlue}{RGB}{0,0,0}        % 描边：黑
\definecolor{MidBlue}{RGB}{70,70,70}      % 次曲线 / 网格：深灰
\tikzset{
  >=Stealth,
  every picture/.append style={font=\small},
  axisline/.style={thick,color=black,-Stealth},
  axislabel/.style={color=black},
  fillsoft/.style={fill=LightRed,fill opacity=0.6,draw=black,thick},
  nodeblue/.style={circle,fill=black!85,draw=black,thick,inner sep=2pt},
  nodewhite/.style={circle,fill=white,draw=black,thick,inner sep=2pt},
  arrblue/.style={->, thick, color=black},
  boxsoft/.style={draw=black, rounded corners=2pt, align=center, inner sep=5pt, fill=black!6},
  % 多曲线 / 多集合区分：以线型 + 标记代替颜色
  curveA/.style={thick,black,solid},
  curveB/.style={thick,black,dashed},
  curveC/.style={thick,black,dotted},
  markA/.style={mark=*,mark size=2.2pt,color=black},
  markB/.style={mark=square*,mark size=2.0pt,color=black},
  markC/.style={mark=triangle*,mark size=2.4pt,color=black},
  % 图内区域 / 点标签（白底避免与图形重叠）
  figlabel/.style={font=\scriptsize,fill=white,inner sep=1pt,outer sep=1pt,text=black},
}
% 图形说明（黑白模式下对图形区域 / 曲线 / 关键点的解释，置于图形下方）
\newcommand{\figexplan}[1]{%
  \par\smallskip\noindent{\small\itshape 图形说明：#1}\par\smallskip}

% ---------- 页眉页脚（书稿格式：左=书名，右=章次与章名）----------
\usepackage{fancyhdr}
% 章标题：居中、大号（书稿章首样式）
\CTEXsetup[format={\centering\LARGE\bfseries}]{chapter}
\CTEXsetup[beforeskip=16pt,afterskip=14pt]{chapter}
\CTEXsetup[format={\large\bfseries}]{section}
\CTEXsetup[format={\normalsize\bfseries}]{subsection}
\pagestyle{fancy}
\fancyhf{}
\providecommand{\booktitle}{微观经济学理论：可行集、偏好与效用}
\providecommand{\lectureno}{第一章}
\providecommand{\lectureshort}{课程讲义}
\fancyhead[L]{\small\itshape \booktitle}
\fancyhead[R]{\small \lectureno\quad \lectureshort}
\fancyfoot[C]{\small \thepage}
\renewcommand{\headrulewidth}{0.4pt}
\fancypagestyle{plain}{\fancyhf{}\fancyfoot[C]{\small \thepage}\renewcommand{\headrulewidth}{0pt}}

% ---------- 正文细节 ----------
\usepackage{setspace}
\setstretch{1.12}
\usepackage{microtype}
\tolerance=900
\emergencystretch=2em
\raggedbottom

% ---------- 超链接与 PDF 书签 ----------
\usepackage{hyperref}
\hypersetup{
  colorlinks=true,
  linkcolor=black,
  urlcolor=black,
  citecolor=black,
  pdfencoding=auto,
  bookmarks=true,
  pdftitle={\booktitle\quad \lectureno：\lectureshort},
  pdfsubject={微观经济学理论书稿}
}
% 章节标题中常见数学符号的安全书签映射（避免 "Improper alphabetic constant"）
\pdfstringdefDisableCommands{%
  \def\leftrightarrow{<->}%
  \def\star{*}%
  \def\S{S}%
  \def\succeq{\text{≽}}%
  \def\succ{\text{≻}}%
  \def\sim{\text{∼}}%
}

% ---------- 常用记号 ----------
\newcommand{\R}{\mathbb{R}}
\newcommand{\N}{\mathbb{N}}
\newcommand{\Z}{\mathbb{Z}}
\newcommand{\Q}{\mathbb{Q}}
\newcommand{\succeqq}{\succeq}
\newcommand{\argmax}{\operatorname*{arg\,max}}
\newcommand{\argmin}{\operatorname*{arg\,min}}
\newcommand{\st}{\text{s.t.}\ }

\setcounter{chapter}{0}
% 缺字回退：主字体缺字符时回退到 Songti SC
\directlua{luaotfload.add_fallback("wbfallback", {"name:Songti SC:mode=node;"})}
\setmainfont{TeX Gyre Termes}[Ligatures=TeX,RawFeature={fallback=wbfallback}]

\begin{document}
\chapter{双产出技术的收益最大化：一个完整可算的数值案例}

\noindent\textbf{本案例的任务。}本案例以一个可以手算到底的两产品企业为对象，
把生产者理论中"转换函数—边际转换率—拉格朗日乘子—收益函数"这条主线
从符号约定走到经济解释，逐步给出全部中间步骤，并对每一处容易出错的符号与方向
做显式辨析。案例所用技术为
\begin{equation*}
F(y_1,y_2)=y_1+\frac{y_2^{2}}{4}-10=0,
\qquad p_1=2,\quad p_2=4,
\end{equation*}
它足够简单（一阶条件是线性的）、又足够丰富（包含机会成本递增、
唯一内点解、可显式写出收益函数、并能在价格比极端时出现角点解），
因此适合作为课堂讲授、作业与考试的标准母题。

\noindent 本案例同时承担三个"辨析"任务，这三处正是初学者最常出错的地方：

\begin{enumerate}
  \item \textbf{符号约定}。转换函数写 $F=0$ 还是 $F\ge 0$，
        可行集在 $F<0$ 一侧还是 $F>0$ 一侧，不是客观事实，而是约定；
        但它决定拉格朗日里写 $-\lambda F$ 还是 $+\lambda F$，决定 $\lambda$ 的符号，
        一旦中途换约定就会全盘出错（\S\ref{sec:convention}）。
  \item \textbf{MRT 的方向}。若以 $MRT_{12}$ 表示"多产 $1$ 单位 $y_1$
        必须放弃的 $y_2$ 数量"，则均衡条件为 $MRT_{12}=p_1/p_2$；
        而当 $MRT_{12}>p_1/p_2$ 时，正确的调整方向是\textbf{减少} $y_1$、
        \textbf{增加} $y_2$，而不是相反（\S\ref{sec:mrt} 给出三种独立验证）。
  \item \textbf{$-\lambda F$ 的性质}。在等式约束问题里，$-\lambda F$ 不是"惩罚项"，
        $\lambda$ 是技术约束的影子值；真正的惩罚项只出现在罚函数法中，
        且惩罚的对象是 $F^{2}$ 而非 $\lambda F$（\S\ref{sec:penalty}）。
\end{enumerate}

\tableofcontents
\clearpage

% ============================================================
%  第一部分：案例设定、收益最大化问题与逐步数值求解
% ============================================================

\section{案例设定：双产出技术与 Shephard 转换函数}
\label{sec:convention}

\subsection{一个"厨房"里的两产品企业}

设想一家面包房，给定一天内可支配的全部要素（一名师傅、一台烤箱、固定数量的面粉与糖），
它可以生产两种产品：

\begin{itemize}
  \item $y_1$：面包的产量（单位：件）；
  \item $y_2$：蛋糕的产量（单位：件）。
\end{itemize}

两种产品都要占用同一批要素，因此之间存在替代关系：多烤一个蛋糕，
就得少烤几个面包。这种"给定要素禀赋下两种产出的技术替代关系"，
在数学上不必回到生产函数去处理，而可以直接用一个\textbf{转换函数}
（transformation function，Shephard 的标准设定）来描述。
本案例的技术设定为
\begin{equation}
F(y_1,y_2)=y_1+\frac{y_2^{2}}{4}-10,
\label{eq:tech}
\end{equation}
市场价格为 $p_1=2$、$p_2=4$（单位：元/件）。目标是选择产出组合最大化收益：
\begin{equation}
\max_{y_1,y_2}\ \pi=p_1y_1+p_2y_2=2y_1+4y_2
\quad\text{s.t.}\quad F(y_1,y_2)=0.
\label{eq:rm}
\end{equation}

在进一步讨论之前，必须先固定 $F$ 的\textbf{符号约定}。
这是整件事的根基：约定一旦确定，$- \lambda F$ 还是 $+\lambda F$、
$\lambda$ 的正负、可行集在哪一侧，就全部锁定了。

\subsection{转换函数的三项约定}

\begin{definition}[转换函数及其约定]
设 $F:\R^{2}_{+}\to\R$ 连续可微，且对每个 $i$ 有 $\partial F/\partial y_i>0$
（产出越多，$F$ 越大，即技术约束是"紧"的、不存在自由处置意义上的浪费被隐藏）。
称 $F$ 为该技术的\textbf{转换函数}，并约定：
\begin{equation}
F(y_1,y_2)
\begin{cases}
=0, & \text{产出组合恰好落在生产可能性边界（PPF）上，技术有效；}\\[2pt]
<0, & \text{产出组合落在 PPF 内部，资源未用尽，技术无效率；}\\[2pt]
>0, & \text{产出组合超出 PPF，技术上不可行。}
\end{cases}
\label{eq:convention}
\end{equation}
\end{definition}

三条约定的含义要逐条落实。

\begin{itemize}
  \item \textbf{为什么 $F=0$ 是边界。}等值线 $F(y)=0$ 是一条曲线，
        它把平面分成两块区域；由于 $F$ 连续，从 $F>0$ 的点走到 $F<0$ 的点
        必定穿过这条曲线，所以它就是集合 $\{y:F(y)\le 0\}$ 的边界，
        也就是生产可能集的上边界。
  \item \textbf{为什么可行的一侧是 $F\le 0$ 而不是 $F\ge 0$。}
        关键在于单调性假设 $\partial F/\partial y_i>0$：
        若某点满足 $F<0$，则把任一产出减少一点点，$F$ 会变得更小，
        于是该点连同它"左下方的所有点"都属于可行集；
        而若 $F>0$，减少产出只能使 $F$ 沿着连续路径下降，
        无论减多少都不可能把它变进 $F\le 0$ 而不丢掉已实现的一部分产出，
        因此 $F>0$ 的点位于要素禀赋所支持的范围之外。
        一句话：$\partial F/\partial y_i>0$ 保证了"内部点可达、外部点不可达"。
  \item \textbf{为什么必须是严格不等号才叫无效率。}
        $F<0$ 表示存在未用尽的要素：可以在不减少 $y_1$ 的条件下增加一些 $y_2$
        （或反之），直到 $F$ 上升到 $0$。因此利润最大化企业绝不会停在严格内部点上。
\end{itemize}

\begin{remark}[约定是任意的，但不能中途更换]
把可行的一侧写成 $F\le 0$ 只是约定。若有人用\ $\widetilde{F}=-F$，
则同一技术的可行集变为 $\widetilde{F}\ge 0$，拉格朗日要写成
$\mathcal{L}=p\cdot y+\mu\,\widetilde{F}$ 才能得到正的乘子。
\textbf{判断一套写法是否自洽，只有一条标准：在最优点上，价格向量必须与所选转换函数的梯度同向且成严格正比例。}
本案例始终遵守定义 \ref{eq:convention}，故拉格朗日写为减去 $\lambda F$，乘子为正。
\end{remark}

除上述三项约定之外还有一层容易被忽略的问题：$F$ 本身的量纲。\eqref{eq:tech} 中的 $\frac14 y_2^{2}$ 看似与 $y_1$ 量纲不合，
严格写法应引入带单位的曲率参数 $\kappa$：
\begin{equation}
F(y_1,y_2)=y_1+\frac{\kappa}{2}y_2^{2}-c,
\qquad [\kappa]=\text{件}^{-1},\quad [c]=\text{件},
\label{eq:tech-general}
\end{equation}
本案例取 $\kappa=\tfrac12$、$c=10$，于是 $\frac{\kappa}{2}y_2^{2}=\frac14 y_2^{2}$，
恰为 \eqref{eq:tech}。此时 $[F]=\text{件}$，
$F_1=1$ 无量纲，$F_2=\kappa y_2$ 也无量纲。
后面 \S\ref{sec:grad} 会用到这一点：$\lambda$ 的单位是元/件，与价格一致。

\subsection{三项约定的数值检验}

取同一条竖线 $y_2=4$ 上的三个点来检验约定 \eqref{eq:convention}：

\begin{equation}
\begin{array}{lll}
\text{点 }(y_1,y_2)=(6,4): & F=6+\tfrac{16}{4}-10=0, & \text{前沿（技术有效）；}\\[2pt]
\text{点 }(y_1,y_2)=(4,4): & F=4+\tfrac{16}{4}-10=-2<0, & \text{内部（要素闲置）；}\\[2pt]
\text{点 }(y_1,y_2)=(8,4): & F=8+\tfrac{16}{4}-10=+2>0, & \text{不可行。}
\end{array}
\label{eq:check-three}
\end{equation}

\begin{figure}[H]
\centering
\begin{tikzpicture}[x=0.95cm,y=0.50cm,line cap=round,line join=round]
  % ---- 可行集（F<=0）阴影 ----
  \fill[black!6] (0,0) -- (0,10)
    -- plot[domain=0:6.3246,samples=70] (\x,{10-\x*\x/4})
    -- (6.3246,0) -- cycle;
  % ---- 坐标轴 ----
  \draw[axisline] (0,0) -- (7.6,0) node[below,axislabel] {$y_2$（蛋糕，件）};
  \draw[axisline] (0,0) -- (0,13.8) node[left,axislabel] {$y_1$（面包，件）};
  % ---- 生产可能性边界 ----
  \draw[curveA,domain=0:6.3246,samples=90] plot (\x,{10-\x*\x/4});
  \node[figlabel,anchor=east] at (6.0,0.55) {$y_1=10-\dfrac{y_2^{2}}{4}$};
  % ---- 检验点所在的竖线 ----
  \draw[curveC] (4,0) -- (4,8.6);
  % ---- 三个检验点 ----
  \draw[nodewhite] (4,6) circle (2.2pt);
  \draw[nodewhite] (4,4) circle (2.2pt);
  \draw[nodewhite] (4,8) circle (2.2pt);
  \fill (4,6) circle (2.2pt);
  \fill (4,4) circle (2.2pt);
  \fill (4,8) circle (2.2pt);
  % ---- 三个标签 ----
  \node[figlabel,anchor=west] at (4.25,8.0) {$(8,4)$：$F=+2>0$，技术不可行};
  \node[figlabel,anchor=west] at (4.25,6.0) {$(6,4)$：$F=0$，生产前沿};
  \node[figlabel,anchor=west] at (4.25,4.0) {$(4,4)$：$F=-2<0$，内部无效率};
  % ---- 区域说明 ----
  \node[figlabel] at (2.2,2.2) {$F<0$：可行域};
  \node[figlabel] at (1.6,11.6) {$F>0$：不可行域};
  % ---- 端点 ----
  \draw (0,0.12) -- (-0.12,0.12) -- (-0.12,0);
  \node[figlabel,anchor=east] at (0.05,10) {$10$};
  \node[figlabel,anchor=north] at (6.3246,0.05) {$\sqrt{40}$};
\end{tikzpicture}
\caption{转换函数的符号约定：粗实线为生产可能性边界 $F(y_1,y_2)=0$，
  其左下方的浅灰区域为可行集 $\{y:F(y)\le 0\}$，右上区域技术不可行。
  同一竖线 $y_2=4$ 上的三个点分别对应 $F>0$、$F=0$、$F<0$ 三种情形。}
\label{fig:convention}
\end{figure}
\figexplan{横轴为 $y_2$（蛋糕），纵轴为 $y_1$（面包）。
粗实线 $y_1=10-y_2^{2}/4$ 是唯一的技术前沿，它向下弯曲，
反映 $y_2$ 越多、再多烤一单位蛋糕要放弃的面包越多（机会成本递增）。
浅灰阴影为可行集，其与两坐标轴之间的部分就是"要素未用尽"的内部点。
虚线竖线穿过三个检验点：最低点 $(4,4)$ 在阴影内（$F=-2$），
中间点 $(6,4)$ 恰在边界上（$F=0$），最高点 $(8,4)$ 在阴影外（$F=+2$），
它在给定要素下无法生产出来。}

\subsection{本案例技术的三条性质}

在动手求解之前，先把技术的形状说清楚，因为二阶条件与唯一性都取决于它。

\begin{enumerate}
  \item \textbf{单调性}。$F_1=\partial F/\partial y_1=1>0$，
        $F_2=\partial F/\partial y_2=\dfrac{y_2}{2}>0$（对 $y_2>0$），
        满足转换函数的单调性要求，故约定 \eqref{eq:convention} 成立。
  \item \textbf{可行集紧（闭且有界）}。由 $F\le 0$ 且 $y\ge 0$ 得
        $0\le y_2\le \sqrt{40}$、$0\le y_1\le 10-\dfrac{y_2^{2}}{4}$，
        可行集是有界闭集；连续函数 $p\cdot y$ 在其上必取到最大值（Weierstrass 定理），
        \textbf{最大值一定存在}，这是后面一切讨论的前提。
  \item \textbf{前沿严格凹、机会成本递增}。把前沿写成显式函数
        \begin{equation}
        y_1=T(y_2)=10-\frac{y_2^{2}}{4},\qquad
        T'(y_2)=-\frac{y_2}{2}<0,\quad T''(y_2)=-\frac12<0,
        \label{eq:ppf-explicit}
        \end{equation}
        $T$ 严格凹，即生产可能集严格凸之外的"边界向外凸"形状：
        随着 $y_2$ 增加，$-(dy_1/dy_2)=y_2/2$ 递增，
        每多烤一单位蛋糕要放弃的面包越来越多。
\end{enumerate}

第三条性质直接给出\textbf{二阶充分条件}：把约束代入目标，
得到只含一个变量的收益函数
\begin{equation}
\pi(y_2)=p_1T(y_2)+p_2y_2
=2\Bigl(10-\frac{y_2^{2}}{4}\Bigr)+4y_2
=20-\frac{y_2^{2}}{2}+4y_2,
\qquad \pi''(y_2)=-1<0,
\label{eq:one-dim}
\end{equation}
它是 $y_2$ 的严格凹函数。因此任何满足一阶条件的内点都是\textbf{唯一的全局最大点}，
不需要再讨论多峰或局部极值问题。这一点在课堂上值得先说，
因为它把"检查二阶条件"变成了看一眼二阶导数符号。

完成平方还能把 \eqref{eq:one-dim} 写成一个更便于口算的形式：
\begin{equation}
\pi(y_2)=28-\frac{(y_2-4)^{2}}{2}.
\label{eq:pi-complete-square}
\end{equation}
沿前沿的收益是一条以 $y_2=4$ 为顶点、最高值为 $28$ 的抛物线。
后面 \S\ref{sec:solve} 用拉格朗日法得到的答案，必须与这条抛物线的顶点完全一致。

\section{收益最大化问题与拉格朗日一阶条件}
\label{sec:foc}

\subsection{问题与几何图像}

把 \eqref{eq:rm} 写成一个正式的最优化问题：
\begin{equation}
\tag{RM}
\max_{y_1,y_2}\ p_1y_1+p_2y_2
\quad\text{s.t.}\quad F(y_1,y_2)=0,\quad y_1\ge 0,\ y_2\ge 0.
\label{eq:rm-formal}
\end{equation}
目标函数 $p\cdot y$ 是线性的，其等值线$p_1y_1+p_2y_2=\bar{\pi}$ 是直线，
斜率（以横轴 $y_2$、纵轴 $y_1$ 度量）为
\begin{equation}
\frac{dy_1}{dy_2}=-\frac{p_2}{p_1}=-\frac{4}{2}=-2.
\label{eq:isor-slope}
\end{equation}
几何上，\eqref{eq:rm-formal} 就是问：在给定斜率的直线族中，
哪一条与前沿 $F=0$ 相交且截距（即收益）最大。
答案是与前沿\textbf{相切}的那一条——这就是后面所有代数要复述的内容。

\subsection{拉格朗日函数与一阶条件}

按定义 \ref{eq:convention} 的约定构造拉格朗日函数：
\begin{equation}
\mathcal{L}(y_1,y_2,\lambda)=p_1y_1+p_2y_2-\lambda F(y_1,y_2).
\label{eq:lagrange}
\end{equation}
对两个产出求偏导并令其为零：
\begin{equation}
\begin{cases}
\dfrac{\partial \mathcal{L}}{\partial y_1}=p_1-\lambda\dfrac{\partial F}{\partial y_1}=0,\\[8pt]
\dfrac{\partial \mathcal{L}}{\partial y_2}=p_2-\lambda\dfrac{\partial F}{\partial y_2}=0,
\end{cases}
\quad\Longleftrightarrow\quad
\begin{cases}
p_1=\lambda F_1,\\
p_2=\lambda F_2.
\end{cases}
\label{eq:foc}
\end{equation}
两式相除消去 $\lambda$（要求 $F_2\neq 0$），得到本案例的核心等式：
\begin{equation}
\boxed{\ \frac{p_1}{p_2}=\frac{F_1}{F_2}\ }
\label{eq:core}
\end{equation}
第三个一阶条件是把 $\mathcal{L}$ 对 $\lambda$ 求偏导，它把约束还回来：
$\partial \mathcal{L}/\partial \lambda=-F(y_1,y_2)=0$，即 $F=0$。

三个未知数 $y_1,y_2,\lambda$，三个方程，可解。
这里只需要一条约束规格：$\nabla F=(F_1,F_2)\neq (0,0)$，
在最优点附近由 $F_1=1$ 自动满足（梯度秩为 $1$）。

\subsection{为什么乘子为正：符号、量纲与经济含义}

由 \eqref{eq:foc} 直接读出：
\begin{equation}
\lambda=\frac{p_1}{F_1}=\frac{p_2}{F_2}.
\label{eq:lambda-ratio}
\end{equation}
本例中 $p_1,p_2>0$ 且 $F_1,F_2>0$，故 $\lambda>0$。这不是巧合：

\begin{itemize}
  \item \textbf{符号方面}。$F$ 沿任一产出增加的方向上升，而价格非负，
        两者之比自然为正。拉格朗日里写 $-\lambda F$ 而不是 $+\lambda F$，
        正是为了保证这个正数出现；若写成 $+\lambda F$，
        解出的 $\lambda$ 会是 $-2$，解释为"价格与技术梯度反向"，徒增混乱。
  \item \textbf{量纲方面}。由 \eqref{eq:tech-general}，
        $[F_1]=1$（无量纲），$[p_1]=$ 元/件，故 $[\lambda]=$ 元/件，
        与价格同量纲。于是 $p=\lambda\nabla F$ 两边量纲一致：
        无单位的梯度方向被乘以一个"元/件"的比例因子，得到价格向量。
  \item \textbf{经济方面}。$\lambda$ 把\textbf{纯技术的方向}
        （$\nabla F$，不含任何价格信息）缩放成\textbf{市场价格的方向}。
        \S\ref{sec:shadow} 将证明它等于收益对技术水平的导数，
        是技术约束的影子值。
\end{itemize}

\section{逐步求解：从梯度到最优产出组合}
\label{sec:solve}

下面按四个步骤把 \eqref{eq:rm-formal} 解到底，每一步的中间结果都单独列出，
便于课堂上逐步板书。

\subsection{第一步：求偏导（梯度分量）}

由 $F=y_1+\dfrac{y_2^{2}}{4}-10$：
\begin{equation}
F_1=\frac{\partial F}{\partial y_1}=1,
\qquad
F_2=\frac{\partial F}{\partial y_2}=\frac{y_2}{2}.
\label{eq:grad-components}
\end{equation}
注意 $F_1$ 是常数，$F_2$ 随 $y_2$ 递增：这条技术对 $y_1$ 而言是"线性的"，
对 $y_2$ 而言呈凸的技术代价，这正是前沿弯曲的来源。

\subsection{第二步：代入一阶条件}

把 \eqref{eq:grad-components} 与价格 $p_1=2,p_2=4$ 代入 \eqref{eq:foc}：
\begin{equation}
2=\lambda\cdot 1,
\qquad
4=\lambda\cdot \frac{y_2}{2}.
\label{eq:foc-numbers}
\end{equation}
第一式立刻给出
\begin{equation}
\lambda=2.
\label{eq:lambda-value}
\end{equation}
代入第二式：$4=2\cdot \dfrac{y_2}{2}=y_2$，故
\begin{equation}
y_2^{*}=4.
\label{eq:y2-star}
\end{equation}

\subsection{第三步：回代前沿求出产出组合与收益}

把 $y_2^{*}=4$ 代入前沿 $F=0$，即显式形式 \eqref{eq:ppf-explicit}：
\begin{equation}
y_1^{*}=10-\frac{(4)^{2}}{4}=10-4=6.
\label{eq:y1-star}
\end{equation}
对应的最优收益（最大可能的总收入）为
\begin{equation}
\pi^{*}=p_1y_1^{*}+p_2y_2^{*}=2\times 6+4\times 4=12+16=28.
\label{eq:pi-star}
\end{equation}
最优解汇总：
\begin{equation}
\boxed{\ (y_1^{*},y_2^{*})=(6,4),\qquad \lambda^{*}=2,\qquad \pi^{*}=28\ }
\label{eq:solution}
\end{equation}

\subsection{第四步：验证核心等式}

核心等式 \eqref{eq:core} 的左边是市场价格比，右边取在最优点 $(6,4)$ 处的技术梯度比：
\begin{equation}
\frac{p_1}{p_2}=\frac{2}{4}=\frac12,
\qquad
\frac{F_1}{F_2}\bigg|_{(6,4)}=\frac{1}{4/2}=\frac{1}{2}.
\label{eq:core-check}
\end{equation}
两边相等，验证通过。同时检查 $p=\lambda\nabla F$：
\begin{equation}
\lambda\nabla F=2\cdot\Bigl(1,\frac{y_2}{2}\Bigr)\bigg|_{y_2=4}
=2\cdot(1,2)=(2,4)=(p_1,p_2)=p.
\label{eq:p-lambda-grad}
\end{equation}
价格向量与技术梯度在最优点上共线，比例因子为 $\lambda=2$。

\subsection{交叉验证一：消元法}

不做拉格朗日，改用"把约束消掉"的老办法。由前沿 $y_1=10-\dfrac{y_2^{2}}{4}$
代入目标，得到单变量问题
\begin{equation}
\max_{0\le y_2\le \sqrt{40}}\ \pi(y_2)=20-\frac{y_2^{2}}{2}+4y_2
=28-\frac{(y_2-4)^{2}}{2}.
\label{eq:onedim}
\end{equation}
一阶条件：$\pi'(y_2)=4-y_2=0\Rightarrow y_2^{*}=4$；
二阶条件：$\pi''=-1<0$，严格凹，故为唯一全局最大点。回代得 $y_1^{*}=6$、$\pi^{*}=28$。
与 \eqref{eq:solution} 完全一致。

两种方法的对应关系值得点明：消元法里的"$4-y_2$"就是"$p_2-\lambda F_2$"
除以某个正因子的结果，二者是同一条边际条件的两种写法；
拉格朗日法的优势在于不必显式解出前沿函数 $T(\cdot)$，
在多产出、多约束的问题里尤其明显。

\subsection{交叉验证二：前沿上的收益对照表}

沿前沿逐点计算收益，可以直接看到最大值出现在何处。
由 \eqref{eq:pi-complete-square}，收益只取决于 $(y_2-4)^{2}$，
表 \ref{tab:frontier-pi} 给出七个代表性组合。

\begin{table}[H]
\centering
\caption{沿生产可能性边界的收益对照（$p_1=2,\ p_2=4$）}
\label{tab:frontier-pi}
\begin{tabular}{ccccc}
\toprule
$y_2$  & $y_1=10-\dfrac{y_2^{2}}{4}$ & $\pi=2y_1+4y_2$ & $28-\dfrac{(y_2-4)^{2}}{2}$ & 位置 \\
\midrule
$0$      & $10.00$  & $20.00$  & $20.00$  & 端点（只产面包）\\
$2$      & $9.00$   & $26.00$  & $26.00$  & 偏离 \\
$3$      & $7.75$   & $27.50$  & $27.50$  & 偏离 \\
$\mathbf{4}$ & $\mathbf{6.00}$ & $\mathbf{28.00}$ & $\mathbf{28.00}$ & \textbf{最优（相切）}\\
$5$      & $3.75$   & $27.50$  & $27.50$  & 偏离 \\
$6$      & $1.00$   & $26.00$  & $26.00$  & 偏离 \\
$\sqrt{40}\approx 6.325$ & $0$ & $25.30$ & $25.30$ & 端点（只产蛋糕）\\
\bottomrule
\end{tabular}
\end{table}

表 \ref{tab:frontier-pi} 说明两件事。第一，最大值确实出现在 $y_2=4$ 而不是别处；
第二，两个端点都明显吃亏：只产面包得到 $20$，只产蛋糕得到 $25.30$，
都低于混合生产的 $28$。也就是说，\textbf{只要价格比不极端，
专业化（走到端点）就不是最优}，这一点在 \S\ref{sec:corner} 会被严格化。

\begin{figure}[H]
\centering
\begin{tikzpicture}[x=0.95cm,y=0.50cm,line cap=round,line join=round]
  % ---- 可行集阴影 ----
  \fill[black!6] (0,0) -- (0,10)
    -- plot[domain=0:6.3246,samples=70] (\x,{10-\x*\x/4})
    -- (6.3246,0) -- cycle;
  % ---- 坐标轴 ----
  \draw[axisline] (0,0) -- (7.6,0) node[below,axislabel] {$y_2$（蛋糕，件）};
  \draw[axisline] (0,0) -- (0,13.8) node[left,axislabel] {$y_1$（面包，件）};
  % ---- 生产前沿 ----
  \draw[curveA,domain=0:6.3246,samples=90] plot (\x,{10-\x*\x/4});
  \node[figlabel,anchor=north east] at (1.15,9.68) {PPF：$F=0$};
  % ---- 等收益线 pi = 20（过端点 (10,0)）----
  \draw[curveB] (0,10) -- (5,0);
  \node[figlabel,anchor=north west] at (0.85,8.3) {$\bar\pi=20$};
  % ---- 等收益线 pi = 25.30（过端点 (0,sqrt40)）----
  \draw[curveC] (0,12.649) -- (6.3246,0);
  \node[figlabel,anchor=south west] at (0.15,12.3) {$\bar\pi\approx 25.30$};
  % ---- 等收益线 pi = 28（与 PPF 相切）----
  \draw[thick,black] (2,10) -- (7,0);
  \node[figlabel,anchor=south west] at (6.0,1.6) {$\bar\pi=28$};
  % ---- 最优 -----端点标记
  \fill (4,6) circle (2.6pt);
  \node[figlabel,anchor=south east] at (3.9,6.15) {$(y_2^{*},y_1^{*})=(4,6)$};
  \draw (0,0.12) -- (-0.12,0.12) -- (-0.12,0);
  \node[figlabel,anchor=east] at (0.05,10) {$10$};
  \node[figlabel,anchor=north] at (6.3246,0.05) {$\sqrt{40}$};
\end{tikzpicture}
\caption{收益最大化的几何：粗实线为生产前沿，
  三条直线为不同收益水平的等收益线。$2y_1+4y_2=28$ 与前沿相切于 $(y_2,y_1)=(4,6)$，
  是所有与前沿有交点的等收益线中位置最高的一条。}
\label{fig:optimal-tangency}
\end{figure}
\figexplan{横轴 $y_2$、纵轴 $y_1$。
粗实线（PPF）向下弯曲；三条等收益线斜率同为 $-p_2/p_1=-2$：
虚线 $\bar\pi=20$ 经过只产面包的端点，
点线 $\bar\pi\approx 25.30$ 经过只产蛋糕的端点，
实线 $\bar\pi=28$ 与前沿恰好相切，黑点即最优点 $(y_2^{*},y_1^{*})=(4,6)$。
在穿过相同两条直线的意义上，$\bar\pi=28$ 的位置"最高"：
任何更高的平行线都会与前沿脱离接触，对应技术做不到的产出组合。}

% ============================================================
%  第二部分：边际转换率、梯度与几何辨析
% ============================================================

\section{边际转换率与均衡调整}
\label{sec:mrt}

\subsection{什么是边际转换率}

"前沿上多生产一单位 $y_1$，必须放弃多少 $y_2$"这个问题由前沿的斜率回答。
对隐函数 $F(y_1,y_2)=0$ 两边全微分：
\begin{equation}
F_1\,dy_1+F_2\,dy_2=0
\quad\Longrightarrow\quad
\frac{dy_2}{dy_1}=-\frac{F_1}{F_2}.
\label{eq:implicit-diff}
\end{equation}

\begin{definition}[边际转换率]
以 $y_1$ 为计价方向的边际转换率定义为
\begin{equation}
MRT_{12}\equiv-\frac{dy_2}{dy_1}=\frac{F_1}{F_2},
\label{eq:mrt}
\end{equation}
读作"多生产 $1$ 单位 $y_1$，技术上必须放弃的 $y_2$ 数量"。
反过来，$MRT_{21}=-dy_1/dy_2=F_2/F_1=1/MRT_{12}$。
\end{definition}

代入本案例的梯度 \eqref{eq:grad-components}：
\begin{equation}
MRT_{12}(y_2)=\frac{F_1}{F_2}=\frac{1}{y_2/2}=\frac{2}{y_2},
\qquad
MRT_{21}(y_2)=\frac{y_2}{2}.
\label{eq:mrt-example}
\end{equation}
于是 $MRT_{12}$ 随 $y_2$ 递减、$MRT_{21}$ 随 $y_2$ 递增：蛋糕烤得越多，
再增加一单位蛋糕要牺牲的面包越多。这再次印证前沿的机会成本递增性质，
也保证了一阶条件给出的驻点是最大值而非最小值。

在最优点 $(6,4)$：
\begin{equation}
MRT_{12}=\frac{2}{4}=\frac12,\qquad MRT_{21}=2.
\label{eq:mrt-at-opt}
\end{equation}

\subsection{均衡条件：技术取舍等于市场交换}

把一阶条件 \eqref{eq:foc} 相除得到的 \eqref{eq:core} 与定义 \eqref{eq:mrt} 合起来，
就是本案例的经济学结论。

\begin{theorem}[双产出收益最大化的内点均衡条件]
设 $p\gg 0$，技术满足 $F_1,F_2>0$，且 $y^*$ 为 \eqref{eq:rm-formal} 的内点解，则
\begin{equation}
MRT_{12}(y^{*})=\frac{p_1}{p_2},
\qquad\text{等价地}\qquad
p=\lambda \nabla F(y^{*}),\ \lambda>0.
\label{eq:equilibrium}
\end{equation}
\end{theorem}

这句话的经济内容值得逐字拆开：

\begin{itemize}
  \item \textbf{左边 $MRT_{12}$ 是技术的边际成本}。要把 $1$ 单位资源从 $y_2$ 挪到 $y_1$，
        技术上要付出 $MRT_{12}$ 单位 $y_2$。
  \item \textbf{右边 $p_1/p_2$ 是市场的边际收益}。挪出来的那 $1$ 单位 $y_1$
        在市场上能卖得 $p_1$，按 $y_2$ 计价就是 $p_1/p_2$ 单位 $y_2$。
  \item \textbf{相等即不再有利可图}。$MRT_{12}=p_1/p_2$ 表示
        "技术上的转换代价"恰好等于"市场上的交换所得"，
        企业无论往哪个方向调整，收益都不再增加。
\end{itemize}

\subsection{三种情形与调整方向（含三份独立验证）}

在尚未到达最优的位置上，比较 $MRT_{12}$ 与 $p_1/p_2$ 即可判断该往哪里调。
判断的可靠办法是看\textbf{收益沿前沿的导数}。把 $y_1$ 作为调整对象，由
\eqref{eq:implicit-diff} 有 $dy_2=-MRT_{12}\,dy_1$，于是
\begin{equation}
\frac{d\pi}{dy_1}=p_1+p_2\frac{dy_2}{dy_1}
=p_1-p_2\,MRT_{12}.
\label{eq:dpi-dy1}
\end{equation}
这条式子把一切都摆明了：多 $1$ 单位 $y_1$ 带来收入 $p_1$，
代价是丢掉 $MRT_{12}$ 单位 $y_2$、值 $p_2\,MRT_{12}$。据此得到表 \ref{tab:mrt-cases}。

\begin{table}[H]
\centering
\caption{$MRT_{12}$ 与 $p_1/p_2$ 的三种情形：本案例的最佳应对}
\label{tab:mrt-cases}
\begin{tabular}{p{3.5cm}p{3.3cm}p{3.9cm}p{2.7cm}}
\toprule
情形 & 经济含义（以 $y_1$ 计） & 本案例中的具体位置 & 最优应对 \\
\midrule
$MRT_{12}>\dfrac{p_1}{p_2}$ & 转产 $y_1$ 的成本高于收益：$p_2MRT_{12}>p_1$，$d\pi/dy_1<0$ & $y_2<4$（例如 $y_2=2$：$MRT_{12}=1>\frac12$） & 减少 $y_1$、\textbf{增加} $y_2$ \\[4pt]
$MRT_{12}=\dfrac{p_1}{p_2}$ & 转换成本恰等于市场所得，$d\pi/dy_1=0$ & $y_2=4$：$MRT_{12}=\frac12$ & 停止调整，达到最优 \\[4pt]
$MRT_{12}<\dfrac{p_1}{p_2}$ & 转产 $y_1$ 的成本低于收益：$p_2MRT_{12}<p_1$，$d\pi/dy_1>0$ & $y_2>4$（例如 $y_2=6$：$MRT_{12}=\frac13<\frac12$） & \textbf{增加} $y_1$、减少 $y_2$ \\
\bottomrule
\end{tabular}
\end{table}

三份互相独立的验证可以消除任何符号疑虑：

\begin{enumerate}
  \item \textbf{边际成本—边际收益验证}（即 \eqref{eq:dpi-dy1}）。
        在 $y_2=2$ 处，$MRT_{12}=1$，多产一个面包的代价是 $1$ 个蛋糕（值 $4$ 元），
        多卖一个面包只多得 $2$ 元，亏 $2$ 元，故\textbf{不要}转产面包。
  \item \textbf{直接求导验证}（把约束代入目标）。由 \eqref{eq:one-dim}
        $\pi'(y_2)=4-y_2$：在 $y_2=2$ 时 $\pi'=2>0$，应增加 $y_2$；
        在 $y_2=6$ 时 $\pi'=-2<0$，应减少 $y_2$。与表 \ref{tab:mrt-cases} 完全一致。
  \item \textbf{解析验证}（完全平方）。由 \eqref{eq:pi-complete-square}，
        $\pi(y_2)=28-\frac12(y_2-4)^2$，最高点在 $y_2=4$；
        任何 $y_2\neq 4$ 的位置都只能通过向 $4$ 靠近来改进。
\end{enumerate}

\begin{remark}[一个高频误判]
常见的错误说法是："若 $MRT_{12}>p_1/p_2$，说明技术上 $y_1$ 更值钱，
所以应该多生产 $y_1$。"方向恰好反了。原因是没有分清\textbf{成本}与\textbf{所得}：
$MRT_{12}$ 是生产 $y_1$ 的边际成本（以 $y_2$ 计），$p_1/p_2$ 是生产 $y_1$ 的边际所得
（同样以 $y_2$ 计）。成本高于所得时应当收缩该项，而不是扩张。
用一句口诀记住：\textbf{"贵的东西少做，便宜的东西多做"——这里的"贵"指的是技术代价。}
\end{remark}

\begin{figure}[H]
\centering
\begin{tikzpicture}[x=0.95cm,y=0.50cm,line cap=round,line join=round]
  % ---- 坐标轴 ----
  \draw[axisline] (0,0) -- (7.6,0) node[below,axislabel] {$y_2$（蛋糕，件）};
  \draw[axisline] (0,0) -- (0,13.8) node[left,axislabel] {$y_1$（面包，件）};
  % ---- 生产前沿 ----
  \draw[curveA,domain=0:6.3246,samples=90] plot (\x,{10-\x*\x/4});
  \node[figlabel,anchor=north east] at (1.15,9.68) {PPF：$F=0$};
  % ---- 等收益线 pi = 26（与前沿交于 (2,9) 与 (6,1)）----
  \draw[curveB] (0,13) -- (6.5,0);
  \node[figlabel,anchor=south west] at (0.15,12.65) {$\bar\pi=26$};
  % ---- 沿前沿的改进方向：左段（从 y2=2 向 y2=3.6）----
  \draw[thick,black,-Stealth,domain=2:3.7,samples=40] plot (\x,{10-\x*\x/4});
  % ---- 沿前沿的改进方向：右段（从 y2=6 向 y2=4.6）----
  \draw[thick,black,-Stealth,domain=6:4.5,samples=40] plot (\x,{10-\x*\x/4});
  % ---- 两个交点 ----
  \fill (4,6) circle (2.6pt);
  \draw[nodewhite] (2,9) circle (2.6pt);
  \draw[nodewhite] (6,1) circle (2.6pt);
  \node[figlabel,anchor=south east] at (3.9,6.15) {最优 $(4,6)$};
  \node[figlabel,anchor=west] at (2.12,9.5) {$A=(2,9)$：\,$MRT_{12}=1>\frac12$};
  \node[figlabel,anchor=west] at (6.12,1.5) {$B=(6,1)$：\,$MRT_{12}=\frac13<\frac12$};
  % ---- 方向说明 ----
  \node[figlabel,anchor=north] at (2.85,7.4) {增加 $y_2$};
  \node[figlabel,anchor=south] at (5.25,2.9) {减少 $y_2$};
  \draw (0,0.12) -- (-0.12,0.12) -- (-0.12,0);
  \node[figlabel,anchor=east] at (0.05,10) {$10$};
  \node[figlabel,anchor=north] at (6.3246,0.05) {$\sqrt{40}$};
\end{tikzpicture}
\caption{$MRT_{12}$ 与价格比的偏离决定调整方向。
  等收益线 $\bar\pi=26$ 与前沿交于 $A=(y_2,y_1)=(2,9)$ 与 $B=(6,1)$，
  两点之间前沿位于该线之上，故从任一端沿前沿向中间移动都能提高收益。}
\label{fig:mrt-adjustment}
\end{figure}
\figexplan{在 $A$ 点，$MRT_{12}=2/y_2=1$，大于价格比 $\frac12$：
多产一单位面包要放弃 $1$ 个蛋糕（值 $4$ 元），而多卖一个面包只得 $2$ 元，
亏本，因此应反向调整——减少面包、增加蛋糕，箭头指向 $y_2$ 增大方向。
在 $B$ 点，$MRT_{12}=1/3$，小于 $\frac12$：多产一单位面包只需放弃 $1/3$ 个蛋糕
（值 $4/3$ 元），却多得 $2$ 元，有利可图，箭头指向 $y_2$ 减小方向。
两个方向的箭头都收敛到中间的实心点 $(y_2,y_1)=(4,6)$，那里 $MRT_{12}=\frac12=p_1/p_2$。}

\section{梯度、法向量与乘子}
\label{sec:grad}

\subsection{为什么梯度是前沿的法向量}

取前沿上过 $y^{*}$ 的任意一条光滑参数曲线 $y(t)$，满足 $F(y(t))\equiv 0$。
对 $t$ 求导：
\begin{equation}
\nabla F(y^{*})\cdot \dot y(0)=F_1\dot y_1+F_2\dot y_2=0.
\label{eq:normal}
\end{equation}
即前沿的任意切向量 $\dot y$ 都与 $\nabla F$ 正交，故 $\nabla F$ 是前沿的\textbf{法向量}。
这是纯几何结论，不含任何价格信息。

一阶条件 \eqref{eq:foc} 则可以重写为向量形式：
\begin{equation}
\boxed{\ p=\lambda\,\nabla F(y^{*}),\qquad \lambda>0\ }
\label{eq:p-lambda-nabla}
\end{equation}
它说的是：在最优点，\textbf{收益目标函数的梯度（就是价格向量 $p$）}
与技术约束的梯度（法向量）\textbf{共线且同向}，比例因子为 $\lambda$。
这正是"相切"的代数表达——两条曲线的切线相同，法向量自然共线。

\subsection{本案例中的三条向量}

在最优点 $(y_1^{*},y_2^{*})=(6,4)$：
\begin{equation}
\nabla F=(F_1,F_2)=\Bigl(1,\frac{4}{2}\Bigr)=(1,2),
\qquad
p=(p_1,p_2)=(2,4),
\qquad
p=2\cdot(1,2)=2\,\nabla F.
\label{eq:vectors}
\end{equation}
于是 $\lambda=2$ 承担了三重身份：

\begin{enumerate}
  \item \textbf{缩放因子}。$p$ 与 $\nabla F$ 同方向，长度是它的 $2$ 倍；
  \item \textbf{分量比率}。由 \eqref{eq:lambda-ratio}，
        $\lambda=p_1/F_1=2/1=2$ 与 $\lambda=p_2/F_2=4/2=2$ 给出同一个数，
        这也是一阶条件相容性的校验；
  \item \textbf{影子值}。$\lambda$ 等于收益对技术水平的导数，见 \S\ref{sec:shadow}。
\end{enumerate}

\begin{figure}[H]
\centering
\begin{tikzpicture}[x=1.15cm,y=1.15cm,line cap=round,line join=round]
  % ---- 受限 viewpoint：局部放大，横纵等尺度 ----
  \begin{scope}
    \clip (2.0,3.7) rectangle (6.45,8.5);
    % 生产前沿
    \draw[curveA,domain=1.9:6.0,samples=90] plot (\x,{10-\x*\x/4});
    % 切线：y1 = 14 - 2 y2
    \draw[curveB] (3.0,8.0) -- (6.0,2.0);
    % 法向量所在的直线（点线，向右上延长）
    \draw[curveC] (4,6) -- (7.4,7.7);
    % 法向量：nabla F 段
    \draw[thick,black,-Stealth] (4,6) -- (5,6.5);
    % 价格向量：整段（长度为 2 倍）
    \draw[very thick,black,-Stealth] (4,6) -- (6,7);
  \end{scope}
  % ---- 坐标轴（局部坐标，不从原点开始）----
  \draw[thin,-Stealth] (2.0,3.7) -- (6.45,3.7) node[below,axislabel] {$y_2$};
  \draw[thin,-Stealth] (2.1,3.7) -- (2.1,8.5) node[left,axislabel] {$y_1$};
  % ---- 标注（在裁剪范围之外绘制，避免被截断）----
  \fill (4,6) circle (2.4pt);
  \node[figlabel,anchor=east] at (3.92,6.02) {切点 $(y_2,y_1)=(4,6)$};
  \node[figlabel,anchor=north west] at (5.05,6.45) {$\nabla F=(1,2)$};
  \node[figlabel,anchor=south west] at (6.05,7.02) {$p=\lambda\nabla F=(2,4)$};
  \node[figlabel,anchor=north west] at (5.05,4.02) {切线 $\dfrac{dy_1}{dy_2}=-2$};
  \node[figlabel,anchor=south east] at (2.6,6.9) {前沿 $F=0$};
\end{tikzpicture}
\caption{梯度、法向量与价格向量的共线关系（局部放大，横纵等尺度）。
  在最优点，价格向量 $p$ 与技术梯度 $\nabla F$ 指向同一方向，
  前者长度恰为后者的 $\lambda=2$ 倍；两者共同垂直于前沿在该点的切线。}
\label{fig:gradient}
\end{figure}
\figexplan{本图比例为 $1:1$，因此"垂直"关系在视觉上成立：
实线箭头从切点出发沿右上方向，短的一段是 $\nabla F=(1,2)$
（分量按 $(y_1,y_2)$ 顺序书写，故在"横轴 $y_2$、纵轴 $y_1$"的图面上表现为向右 $2$、向上 $1$），
长的一段是 $p=\lambda\nabla F=(2,4)$，长度恰为两倍。
虚线是前沿的切线，斜率为 $dy_1/dy_2=-2$；点线是法线延长线，
其斜率为 $1/2$，与切线斜率互为负倒数。需注意：本图纵轴不从原点开始，属于局部放大。}

\subsection{量纲与直觉的再次确认}

由 \S\ref{sec:convention} 的量纲约定，$[\nabla F]=1$（无量纲），
$[p]=$ 元/件，故 $[\lambda]=$ 元/件。这让 \eqref{eq:p-lambda-nabla} 读起来非常自然：
\begin{equation*}
\underbrace{p}_{\text{元/件}}
=\underbrace{\lambda}_{\text{元/件}}\cdot
\underbrace{\nabla F}_{\text{无单位的技术方向}} .
\end{equation*}
技术决定"往哪个方向生产更有效"（梯度），市场决定"哪个方向更值钱"（价格），
$\lambda$ 是两者之间的汇率。当所有价格按同一比例变动时，
$p$ 的长度变了而方向不变，\eqref{eq:core} 只涉及价格比，
因此\textbf{最优产出组合不变}——这正是收益最大化行为的零次齐次性，
与消费者理论中马歇尔需求对 $(p,w)$ 零次齐次的性质完全对应。

\section{几何辨析：切线斜率、MRT 与梯度比值}
\label{sec:geometry-check}

同一个几何事实可以有四五种写法，混淆就发生在这些写法之间。
表 \ref{tab:geometry} 把它们并列摆出，并给出本案例在最优点 $(6,4)$ 处的具体数值。

\begin{table}[H]
\centering
\caption{前沿几何的五种写法及其在最优点 $(y_1,y_2)=(6,4)$ 处的取值}
\label{tab:geometry}
\begin{tabular}{p{4.0cm}p{3.3cm}p{2.2cm}p{3.8cm}}
\toprule
对象 & 一般表达式 & 本例取值 & 含义与读法 \\
\midrule
前沿切线斜率 & $\dfrac{dy_1}{dy_2}=-\dfrac{F_2}{F_1}$ & $-2$ & 图上看得见的倾角；纵轴量除以横轴量 \\[4pt]
$MRT_{21}$ & $-\dfrac{dy_1}{dy_2}=\dfrac{F_2}{F_1}$ & $2$ & 多产 $1$ 单位 $y_2$ 要放弃的 $y_1$ \\[4pt]
$MRT_{12}$ & $-\dfrac{dy_2}{dy_1}=\dfrac{F_1}{F_2}$ & $\dfrac12$ & 多产 $1$ 单位 $y_1$ 要放弃的 $y_2$ \\[4pt]
梯度分量之比 & $\dfrac{F_1}{F_2}=\dfrac{p_1}{p_2}$ & $\dfrac12$ & 均衡条件本身，与前一行同一个数 \\[4pt]
梯度向量本身 & $(F_1,F_2)=(1,2)$ & $(1,2)$ & 按 $(y_1,y_2)$ 顺序书写的\textbf{法}向量 \\
\bottomrule
\end{tabular}
\end{table}

三条容易记混的经验规则：

\begin{enumerate}
  \item \textbf{带负号的那一行一定是"应运而生"的量}。$MRT$ 定义为"放弃量"，
        恒为非负；而斜率可以是负的。写到斜率时记得带负号，
        写到 $MRT$ 时它是斜率的相反数（或倒数，取决于方向）。
  \item \textbf{$MRT_{ij}$ 的下标顺序决定用哪个偏导在分子}。
        $MRT_{ij}=F_i/F_j$（$i$ 在分子，$j$ 在分母）。本例 $MRT_{12}=F_1/F_2=1/2$，
        $MRT_{21}=F_2/F_1=2$。
  \item \textbf{梯度是法向量，与切线无关}。写 $(F_1,F_2)$ 时按 $(y_1,y_2)$ 的顺序，
        画到"横轴 $y_2$、纵轴 $y_1$"的图面上要\textbf{交换分量}得到 $(F_2,F_1)$；
        这一步漏掉的人，画出来的箭头会指向切线的镜像方向。
\end{enumerate}

% ============================================================
%  第三部分：乘子的经济含义与一个必须纠正的说法
% ============================================================

\section{乘子的经济含义：影子值与比较静态}
\label{sec:shadow}

\subsection{把技术水平写成参数}

要解释 $\lambda$，必须让技术约束可以"放松一点点"。把 \eqref{eq:tech} 改写为
\begin{equation}
F(y_1,y_2;c)=y_1+\frac{y_2^{2}}{4}-c=0,
\label{eq:tech-param}
\end{equation}
其中 $c>0$ 是技术水平的标量指标：$c$ 越大，同样的 $y_2$ 能挤出越多的 $y_1$，
前沿整体向东北方向扩张。本案例 $c=10$。对应的最大收益记为
\begin{equation}
R(p_1,p_2;c)\equiv \max_{y}\ \{p_1y_1+p_2y_2:\ F(y;c)=0\},
\label{eq:revenue-function}
\end{equation}
即生产者的\textbf{收益函数}：给定价格与技术，企业所能达到的最大收入。

\subsection{一般解：把 $p$ 也留着}

现在把价格也留成符号，一次性解出全部结果。一阶条件仍是
$p_1=\lambda F_1=\lambda$、$p_2=\lambda F_2=\lambda y_2/2$，故
\begin{equation}
\lambda=p_1,
\qquad
\frac{y_2}{2}=\frac{p_2}{\lambda}=\frac{p_2}{p_1}
\ \Longrightarrow\
y_2^{*}(p;c)=\frac{2p_2}{p_1}.
\label{eq:gen-y2}
\end{equation}
回代前沿 $y_1=c-y_2^{2}/4$：
\begin{equation}
y_1^{*}(p;c)=c-\frac{1}{4}\Bigl(\frac{2p_2}{p_1}\Bigr)^{2}
=c-\Bigl(\frac{p_2}{p_1}\Bigr)^{2}.
\label{eq:gen-y1}
\end{equation}
代入目标：
\begin{align}
R(p_1,p_2;c)
&=p_1\Bigl[c-\Bigl(\frac{p_2}{p_1}\Bigr)^{2}\Bigr]
+p_2\cdot\frac{2p_2}{p_1}
=p_1c-\frac{p_2^{2}}{p_1}+\frac{2p_2^{2}}{p_1} \notag\\
&=p_1c+\frac{p_2^{2}}{p_1}.
\label{eq:gen-R}
\end{align}
代入本案例 $p_1=2,p_2=4,c=10$：
\begin{equation}
\lambda=p_1=2,\qquad
y_2^{*}=\frac{2\times 4}{2}=4,\qquad
y_1^{*}=10-\Bigl(\frac{4}{2}\Bigr)^{2}=10-4=6,\qquad
R=2\times 10+\frac{16}{2}=28,
\label{eq:gen-check}
\end{equation}
与逐步求解的 \eqref{eq:solution} 完全一致。
注意 \eqref{eq:gen-R} 只在 $y_1^{*}\ge 0$ 即 $p_2/p_1\le \sqrt{c}$ 时有效，
越界后要换成角点解，见 \S\ref{sec:corner}。

\subsection{三条偏导数与两条经典定理}

对 \eqref{eq:gen-R} 直接求偏导：
\begin{equation}
\frac{\partial R}{\partial c}=p_1,
\qquad
\frac{\partial R}{\partial p_1}=c-\Bigl(\frac{p_2}{p_1}\Bigr)^{2}=y_1^{*},
\qquad
\frac{\partial R}{\partial p_2}=\frac{2p_2}{p_1}=y_2^{*}.
\label{eq:R-partials}
\end{equation}
三条式子各自对应一条一般原理：

\begin{enumerate}
  \item \textbf{乘子即影子值}。由包络定理，
        在最优解处对参数求导只需对拉格朗日函数求偏导：
        \begin{equation}
        \frac{\partial R}{\partial c}
        =\frac{\partial \mathcal{L}}{\partial c}\bigg|_{y=y^{*},\lambda=\lambda^{*}}
        =-\lambda\,\frac{\partial F}{\partial c}
        =-\lambda(-1)=\lambda.
        \label{eq:envelope}
        \end{equation}
        而由 \eqref{eq:gen-R}，$\partial R/\partial c=p_1$，两者一致：
        $\lambda=p_1=2$。这里 $-\partial F/\partial c=1$ 是"参数放松一单位的净额"。
  \item \textbf{Hotelling 型引理}。收益函数对产出价格的偏导等于该产品的最优产量：
        $\partial R/\partial p_i=y_i^{*}$。本例 $\partial R/\partial p_1=10-4=6=y_1^{*}$，
        $\partial R/\partial p_2=4=y_2^{*}$，与之一一对应。
  \item \textbf{齐次性与凸性}。$R$ 对 $(p_1,p_2)$ 一次齐次：
        $R(tp_1,tp_2;c)=tp_1c+\frac{t^{2}p_2^{2}}{tp_1}=tR(p;c)$，
        与"所有价格同比例变动不改变最优产出组合"互为表里。
        对 $p$ 凸：由 \eqref{eq:gen-R}，
        \begin{equation}
        D^{2}_{p}R=
        \begin{pmatrix}
        \dfrac{2p_2^{2}}{p_1^{3}} & -\dfrac{2p_2}{p_1^{2}}\\[6pt]
        -\dfrac{2p_2}{p_1^{2}} & \dfrac{2}{p_1}
        \end{pmatrix},
        \qquad
        \det D^{2}_{p}R=\frac{4p_2^{2}}{p_1^{4}}-\frac{4p_2^{2}}{p_1^{4}}=0,
        \label{eq:hessian-R}
        \end{equation}
        主子式非负，故半正定；此处退化为秩 $1$，
        原因是本例技术只有一条一维前沿（没有任何"平段"），
        供给对价格的反应被限制在这条曲线上。
\end{enumerate}

\subsection{把 \texorpdfstring{$\lambda=2$}{lambda=2} 翻译成一句话}

由 \eqref{eq:envelope}，$\lambda=\partial R/\partial c=2$ 的含义是：
\textbf{技术水平放松一单位，最大收益上升 $2$ 元。}
具体到本案例还有一个更直观的解释：在这条技术上，
$c$ 每增加 $1$，前沿是沿着 $y_1$ 轴方向向外平移 $1$ 个单位
（因为 $c$ 只出现在 $y_1=c-y_2^{2}/4$ 的截距里），
而 $y_1$ 的边际收益就是它的价格 $p_1=2$，故收益增量恰为 $2$ 元。
数值检验：把 $c$ 由 $10$ 调到 $10.01$，由 \eqref{eq:gen-R}
\begin{equation}
R(2,4;10.01)=2\times 10.01+\frac{16}{2}=20.02+8=28.02,
\qquad
\Delta R=0.02=\lambda\times 0.01.
\label{eq:lambda-numeric}
\end{equation}
注意这是\textbf{一阶近似恰好精确}的例子：$R$ 关于 $c$ 是线性的，
所以不只是"近似增加 $0.02$"，而是精确增加 $0.02$。

\begin{figure}[H]
\centering
\begin{tikzpicture}[x=1.10cm,y=0.55cm,line cap=round,line join=round]
  % ---- 坐标轴 ----
  \draw[axisline] (7.6,0) -- (12.5,0) node[below,axislabel] {技术参数 $c$};
  \draw[axisline] (8,23.2) -- (8,32.5) node[left,axislabel] {最大收益 $R(p;c)$};
  % ---- R(c) 直线 ----
  \draw[curveA] (8,24) -- (12,32);
  % ---- 点 (10,28) ----
  \fill (10,28) circle (2.6pt);
  \node[figlabel,anchor=south east] at (9.95,28.1) {$(10,\,28)$};
  % ---- 斜率三角形 ----
  \draw[curveC] (10,28) -- (11,28) -- (11,30);
  \node[figlabel,anchor=north] at (10.5,27.9) {$\Delta c=1$};
  \node[figlabel,anchor=west] at (11.05,29) {$\Delta R=\lambda=2$};
  % ---- 直线标签 ----
  \node[figlabel,anchor=south west,rotate=63] at (11.3,29.4) {$R=2c+8$};
  \node[figlabel,anchor=north west] at (8.15,31.6) {斜率 $=\dfrac{\partial R}{\partial c}=\lambda$};
  % ---- 刻度 ----
  \node[figlabel,anchor=north] at (10,22.9) {$10$};
  \node[figlabel,anchor=east] at (7.95,28) {$28$};
\end{tikzpicture}
\caption{收益函数关于技术参数的比较静态：
$R(2,4;c)=2c+8$ 是一条斜率为 $\lambda=2$ 的直线。
技术水平每放松一单位，最大收益增加 $2$ 元。}
\label{fig:lambda-shadow}
\end{figure}
\figexplan{横轴为技术水平参数 $c$，纵轴为给定价格下的最大收益。
直线在 $c=10$ 处的高度是 $28$，与逐步求解的答案一致；
由 $(10,28)$ 向右走一单位到 $(11,30)$，高度正好上升 $2$，
这个比值 $2/1$ 就是 $\partial R/\partial c$，也等于拉格朗日乘子 $\lambda$。
因为是直线，把 $c$ 调到 $10.01$ 时收益精确增加 $0.02$，没有近似误差。}

\section{一个必须纠正的说法：\texorpdfstring{$-\lambda F$}{-lambda F} 不是"惩罚项"}
\label{sec:penalty}

在口头讲授中，流行这样一种说法：

\begin{quote}
\textsl{"拉格朗日里的 $-\lambda F$ 是惩罚项：$F>0$ 不可行时它被压低，
$F<0$ 无效率时也被压低，只有 $F=0$ 时惩罚消失，所以优化器自动被推到前沿上。"}
\end{quote}

这段话的\textbf{前半句对、后半句错}，而且在最关键的环节上符号写反了。
逐条澄清如下。

\begin{enumerate}
  \item \textbf{$F>0$ 的一侧确实"被压低"，但这不重要}。
        在 $\lambda>0$ 且 $F>0$ 时，$-\lambda F<0$，拉格朗日值小于目标值。
        问题是：$F>0$ 的点根本不在可行集里，最大化从不去那里比较，
        这一侧的"压低"没有参与任何计算。
  \item \textbf{$F<0$ 的一侧恰恰是"被抬高"，不是"被压低"}。
        $F<0$ 时 $-\lambda F>0$，拉格朗日值\textbf{大于}目标值。
        这可以立刻用数字验证：取前沿内部的点 $(y_1,y_2)=(4,4)$，
        $\mathcal{L}=2\times 4+4\times 4-2\times(-2)=32+4=36>24=\pi(4,4)$，
        而前沿上的最优收益仅为 $28$，可见拉格朗日在可行域之外的值没有经济意义。
        若真按"惩罚"理解，一个效率低下的点反倒应当被\textbf{奖励} $4$，这显然说不通。
  \item \textbf{等式约束问题里根本不存在"惩罚"这一步}。
        最大化是在满足 $F=0$ 的集合上进行的，
        第三个一阶条件 $\partial\mathcal{L}/\partial\lambda=0$ 已经把搜索限制在曲面上；
        在曲面上 $-\lambda F\equiv 0$，所谓的惩罚项恒等于零。
\end{enumerate}

\begin{table}[H]
\centering
\caption{流行说法与正确事实的逐条对照}
\label{tab:penalty}
\begin{tabular}{p{4.6cm}p{4.6cm}p{4.6cm}}
\toprule
流行说法 & 正确事实 & 依据 \\
\midrule
"$F<0$ 时 $-\lambda F$ 把目标压低" & $F<0$ 时 $-\lambda F>0$，目标被\textbf{抬高} &
代入 $F=-2,\lambda=2$ \\[4pt]
"惩罚让最优解落到前沿上" & 是第三个一阶条件（约束本身）把解限制在前沿上 &
$\partial\mathcal{L}/\partial\lambda=-F=0$ \\[4pt]
"$F>0$ 会被惩罚，所以不会跑到外面" & $F>0$ 不在可行集，与惩罚无关 &
可行集定义为 $\{y:F(y)\le 0\}$ \\[4pt]
"$\lambda$ 是惩罚系数" & $\lambda$ 是技术约束的影子值 $\partial R/\partial c$ &
式 \eqref{eq:envelope} 与本例 $\lambda=2$ 元/件 \\
\bottomrule
\end{tabular}
\end{table}

\subsection{若一定要谈"惩罚"，正确写法是这样的}

真正带惩罚的最优化是\textbf{罚函数法}：把等式约束去掉，改为最大化
\begin{equation}
\max_{y}\ \Phi_{\rho}(y)=p\cdot y-\frac{\rho}{2}\bigl[F(y)\bigr]^{2},
\qquad \rho>0,\ \rho\to\infty.
\label{eq:penalty-method}
\end{equation}
这里被惩罚的是 $F^{2}$，对 $F>0$ 与 $F<0$ \textbf{对称地}惩罚，
且在 $F=0$ 处惩罚为零、梯度也为零；$\rho$ 越大，解越贴近前沿。
它与拉格朗日法的区别在于两点：一是 $-\frac{\rho}{2}F^{2}$ 是二次的、对称的，
而 $-\lambda F$ 是一次、反对称的；二是 $\rho$ 不是变量而是趋于无穷的参数，
而 $\lambda$ 是与解同时求出来的有限乘子。两者在极限下才相通：
$\rho F(y_{\rho})\to \lambda$。

\begin{figure}[H]
\centering
\begin{tikzpicture}[x=2.60cm,y=1.05cm,line cap=round,line join=round]
  % ---- 坐标轴 ----
  \draw[axisline] (-1.35,0) -- (1.45,0) node[below,axislabel] {$F(y)$};
  \draw[axisline] (0,-2.85) -- (0,2.85) node[left,axislabel] {附加项取值};
  % ---- 区域底色：F<0 可行，F>0 不可行 ----
  \fill[black!5] (-1.3,-2.8) rectangle (0,2.8);
  \node[figlabel] at (-0.72,2.45) {$F<0$：可行（内部）};
  \node[figlabel] at (0.72,2.45) {$F>0$：不可行};
  % ---- 直线 -lambda F ----
  \draw[curveA] (-1.2,2.4) -- (1.2,-2.4);
  \node[figlabel,anchor=south east] at (-0.62,1.3) {$-\lambda F$（$\lambda=2$）};
  % ---- 抛物线 -(rho/2) F^2 ----
  \draw[curveB,domain=-1.3:1.3,samples=80] plot (\x,{-\x*\x});
  \node[figlabel,anchor=north] at (0.72,-0.72) {$-\dfrac{\rho}{2}F^{2}$（$\rho=2$）};
  % ---- F=0 处的关键点 ----
  \fill (0,0) circle (2.4pt);
  \node[figlabel,anchor=north west] at (0.06,0.06) {$F=0$：两项都归零};
  % ---- 取值对比标注 ----
  \node[figlabel,anchor=north west] at (-1.28,2.28) {$+4$（反而抬高）};
  \node[figlabel,anchor=south west] at (-1.28,-1.66) {$-1.69$（压低）};
\end{tikzpicture}
\caption{$-\lambda F$ 与真正的惩罚项 $-\dfrac{\rho}{2}F^{2}$ 的形状对比。
前者在可行的一侧（$F<0$）取值为正，是"抬高"而非"压低"；
后者在 $F=0$ 两侧对称下降，才是惩罚。}
\label{fig:penalty}
\end{figure}
\figexplan{横轴为偏离量 $F$，纵轴为附加到目标上的项。
实直线 $-\lambda F$ 过原点向下倾斜：$F>0$ 时取负值，$F<0$ 时取正值 $+4$
（对应 $F=-2,\lambda=2$），因此它对"效率低下的内部点"是\textbf{奖励}而非惩罚，
这与"把解推到前沿"的直觉恰好相反。
虚线抛物线 $-\frac{\rho}{2}F^{2}$ 以 $F=0$ 为顶点，向两侧对称下降，
在 $F=\pm1.3$ 处取 $-1.69$，这才是惩罚函数的形状。
两条曲线在 $F=0$ 处同为零，这也是它们唯一的交点有意义的原因。}

\subsection{把约束写成不等式，才是严格的说法}

还有一种把直觉救回来的办法：干脆把技术约束写成不等式形式
\begin{equation}
\max_{y}\ p\cdot y
\quad\text{s.t.}\quad F(y)\le 0,\quad y\ge 0,
\label{eq:inequality-form}
\end{equation}
即用生产可能集 $\{y:F(y)\le 0\}$ 而不是它的边界来描述技术。
其 Kuhn--Tucker 条件为：存在 $\mu\ge 0$ 使
\begin{equation}
p=\mu\,\nabla F(y^{*}),\qquad
F(y^{*})\le 0,\qquad
\mu\,F(y^{*})=0.
\label{eq:kt}
\end{equation}
现在"为什么最优解一定在前沿上"就有了一个干净的论证：

\begin{itemize}
  \item 若 $y^{*}$ 是\textbf{严格内部点}，$F(y^{*})<0$，
        由互补松弛 $\mu F=0$ 得 $\mu=0$，代入第一条得 $p=0$，
        与 $p=(2,4)\gg 0$ 矛盾。故严格内部点不可能是最优。
  \item 若 $\mu>0$，则由互补松弛必有 $F(y^{*})=0$，即最优解\textbf{恰在边界（前沿）上}。
  \item 本例解出 $\mu=\lambda=2>0$，落在第二种情形，与前面的计算一致。
\end{itemize}

结论是：\textbf{把解推到前沿上的机制是互补松弛条件，不是"惩罚"。}
 $-\lambda F$ 的正确身份是"目标梯度与约束梯度的配平项"，
它的系数 $\lambda$ 是影子值，而不是罚系数。这段论证也与
马歇尔需求求解中的库恩--塔克四条件一脉相承：
可行性、驻点、乘子非负、互补松弛，四条缺一不可。

% ============================================================
%  第四部分：角点解、对偶、课堂讲授版与小结
% ============================================================

\section{拓展一：什么时候出现角点解}
\label{sec:corner}

\subsection{内点解的有效区间}

一般解 \eqref{eq:gen-y1} 给出的 $y_1^{*}=c-(p_2/p_1)^{2}$ 必须非负，否则它描述的
"内点"落在第一象限之外、根本不在前沿上。于是内点解的有效条件是
\begin{equation}
\Bigl(\frac{p_2}{p_1}\Bigr)^{2}\le c
\quad\Longleftrightarrow\quad
\frac{p_2}{p_1}\le \sqrt{c}=\sqrt{10}\approx 3.1623.
\label{eq:interior-condition}
\end{equation}
一旦价格比越过这个门槛，最优点就撞到端点 $y_1=0$，产生\textbf{角点解}
（完全专业化于蛋糕）。完整的解可以写成一个统一的表达式：
\begin{equation}
y_2^{*}(p;c)=\min\Bigl\{\frac{2p_2}{p_1},\ 2\sqrt{c}\Bigr\},
\qquad
y_1^{*}(p;c)=c-\frac{\bigl(y_2^{*}\bigr)^{2}}{4}.
\label{eq:corner-solution}
\end{equation}
由于本案例的前沿只在第一象限的一小段上有定义
（$0\le y_2\le 2\sqrt{c}$），
$2\sqrt{c}=\sqrt{40}\approx 6.3246$ 正好是端点。

几何上，门槛 \eqref{eq:interior-condition} 的来源一目了然。
前沿斜率 $dy_1/dy_2=-y_2/2$ 在 $0\le y_2\le 2\sqrt{c}$ 上的取值范围是
$[-\sqrt{c},\,0]$；等收益线的斜率是 $-p_2/p_1$。
只要 $-p_2/p_1$ 落在这个区间之外，直线就只能在端点处"靠住"前沿，
切点不复存在。$-p_2/p_1<-\sqrt{c}$ 即 $p_2/p_1>\sqrt{c}$，与 \eqref{eq:interior-condition} 一致。

\subsection{价格比变动的比较静态}

固定 $p_1=2$、$c=10$，让 $p_2$ 变化，用 \eqref{eq:corner-solution} 逐项计算，
结果见表 \ref{tab:corner}。

\begin{table}[H]
\centering
\caption{价格比与最优产出组合（固定 $p_1=2$、$c=10$，门槛 $p_2/p_1=\sqrt{10}\approx 3.162$）}
\label{tab:corner}
\begin{tabular}{ccccccc}
\toprule
$p_2$ & $p_2/p_1$ & 判定 & $y_2^{*}$ & $y_1^{*}$ & $\pi^{*}$ & 位置 \\
\midrule
$1$      & $0.50$  & 内点 & $1.00$   & $9.75$   & $20.50$  & 以面包为主 \\
$2$      & $1.00$  & 内点 & $2.00$   & $9.00$   & $22.00$  & 以面包为主 \\
$4$      & $2.00$  & 内点 & $\mathbf{4.00}$ & $\mathbf{6.00}$ & $\mathbf{28.00}$ & \textbf{本案例}\\
$6$      & $3.00$  & 内点 & $6.00$   & $1.00$   & $38.00$  & 以蛋糕为主 \\
$2\sqrt{10}\approx 6.325$ & $3.162$ & 临界 & $6.325$  & $0$      & $40.00$  & 恰好专业化 \\
$8$      & $4.00$  & 角点 & $6.325$  & $0$      & $50.60$  & 完全专业化于蛋糕 \\
$16$     & $8.00$  & 角点 & $6.325$  & $0$      & $101.19$ & 完全专业化于蛋糕 \\
\bottomrule
\end{tabular}
\end{table}

表 \ref{tab:corner} 读出三条规律。其一，$y_2^{*}$ 对价格比的反应在内点区间内是线性的
（$y_2^{*}=2p_2/p_1$），超过门槛后完全"冻结"在 $2\sqrt{c}$，再也动不了：
这就是\textbf{供给曲线在角点处的弯折}。其二，$y_1^{*}$ 随价格比上升而单调下降，
到达门槛后归零并保持为零。其三，收益随价格单调上升，且在门槛处仍是连续的
（$p_2=2\sqrt{10}$ 时，内点公式给 $p_1c+p_2^{2}/p_1=20+40/2=40$，
角点公式给 $p_2\cdot 2\sqrt{c}=2\sqrt{10}\cdot 2\sqrt{10}=40$，两者相同），
但\textbf{导数发生了跳跃}：产出的调不动了，收入还能继续涨。

\begin{figure}[H]
\centering
\begin{tikzpicture}[x=0.72cm,y=0.38cm,line cap=round,line join=round]
  % ================= 左图：p2/p1=8，角点（专业化于蛋糕） =================
  \begin{scope}
    \fill[black!6] (0,0) -- (0,10)
      -- plot[domain=0:6.3246,samples=60] (\x,{10-\x*\x/4})
      -- (6.3246,0) -- cycle;
    \draw[axisline] (0,0) -- (7.4,0) node[below,axislabel] {$y_2$};
    \draw[axisline] (0,0) -- (0,13.6) node[left,axislabel] {$y_1$};
    \draw[curveA,domain=0:6.3246,samples=80] plot (\x,{10-\x*\x/4});
    % 等收益线：斜率 -8，过点 (6.3246,0)
    \draw[curveB] (4.825,12) -- (6.3246,0);
    \node[figlabel,anchor=south west] at (4.55,11.2) {斜率 $-8$};
    \fill (6.3246,0) circle (2.6pt);
    \node[figlabel,anchor=north west] at (5.4,1.2) {角点 $(y_1,y_2)=(0,\sqrt{40})$};
    \node[figlabel,anchor=south west] at (0.35,12.6) {(a) $p_2/p_1=8>\sqrt{10}$};
    \node[figlabel,anchor=east] at (0.05,10) {$10$};
  \end{scope}
  % ================= 右图：p2/p1=0.25，靠近另一端 =================
  \begin{scope}[xshift=8.2cm]
    \fill[black!6] (0,0) -- (0,10)
      -- plot[domain=0:6.3246,samples=60] (\x,{10-\x*\x/4})
      -- (6.3246,0) -- cycle;
    \draw[axisline] (0,0) -- (7.4,0) node[below,axislabel] {$y_2$};
    \draw[axisline] (0,0) -- (0,13.6) node[left,axislabel] {$y_1$};
    \draw[curveA,domain=0:6.3246,samples=80] plot (\x,{10-\x*\x/4});
    % 等收益线：斜率 -0.25，过点 (0.5, 9.9375)
    \draw[curveB] (0,10.0625) -- (6.3246,8.4813);
    \node[figlabel,anchor=south west] at (2.4,9.4) {斜率 $-0.25$};
    \fill (0.5,9.9375) circle (2.6pt);
    \node[figlabel,anchor=south west] at (0.75,10.6) {$(y_1,y_2)=(9.94,\,0.5)$};
    \node[figlabel,anchor=south west] at (0.35,12.6) {(b) $p_2/p_1=0.25<\sqrt{10}$};
    \node[figlabel,anchor=east] at (0.05,10) {$10$};
  \end{scope}
\end{tikzpicture}
\caption{价格比的两个极端：左图蛋糕相对价格极高，直线比前沿最陡处还陡，
  最优只能停在 $y_1=0$ 的端点（角点解）；右图蛋糕相对价格极低，
  直线近乎水平，最优靠近 $y_2=0$ 一端，但仍是真正的切点（内点解）。}
\label{fig:corner}
\end{figure}
\figexplan{左图虚线的斜率为 $-8$，绝对值大于前沿端点处的斜率 $-3.162$，
它无法在前沿内部找到斜率相同的位置，只能"顶"在右下端点，角点由此产生。
右图虚线斜率为 $-0.25$，落在前沿斜率取值范围之内，与前沿在
$(y_2,y_1)=(0.5,9.94)$ 处真正相切，属于内点解。
两图合起来说明：本例的两个边界，$y_1=0$ 一侧会被价格逼出角点解，
而 $y_2=0$ 一侧只在 $p_2=0$ 的极限下才到达，正常运行中都是内点。}

\section{拓展二：收益函数与教科书对偶体系的位置}
\label{sec:duality}

\subsection{收益函数的四条性质}

把 §\S\ref{sec:shadow} 中得到的收益函数一般化。对任意技术 $F$，
定义 $R(p)=\max\{p\cdot y:F(y)=0\}$，则在适当的光滑性与约束规格下：

\begin{enumerate}
  \item \textbf{对价格一次齐次}：$R(tp)=tR(p)$，$t>0$；
  \item \textbf{对价格单调非减}：$p'\ge p\Rightarrow R(p')\ge R(p)$；
  \item \textbf{对价格凸}：$R$ 是线性函数的上确界，故必凸；
  \item \textbf{Hotelling 引理}：$\nabla_p R(p)=y^{*}(p)$（若最优点唯一），
        本案例已由 \eqref{eq:R-partials} 直接验证。
\end{enumerate}
本案例的显式表达式为（内点区间内）
\begin{equation}
R(p_1,p_2;c)=p_1c+\frac{p_2^{2}}{p_1},
\qquad
R(p;c)=\max\Bigl\{p_1c+\frac{p_2^{2}}{p_1},\ 2p_2\sqrt{c}\Bigr\}\ \text{（含角点情形）}.
\label{eq:R-full}
\end{equation}

\subsection{与消费者理论的对偶映射}

第五讲建立的消费者理论对偶框架，在生产侧有一一对应的版本，
表 \ref{tab:duality} 是二者的对照，便于把两个板块串起来复习。

\begin{table}[H]
\centering
\caption{消费者理论与生产者理论的对象对照}
\label{tab:duality}
\begin{tabular}{p{4.4cm}p{4.4cm}p{4.6cm}}
\toprule
消费者侧 & 生产者侧 & 备注 \\
\midrule
效用函数 $u(x)$ & 转换函数 $F(y)$（或生产可能集 $Y$） & 描述"什么是可行的/如何排序" \\[4pt]
间接效用 $v(p,w)=\max u\ \text{s.t.}\ p\cdot x\le w$ & 收益函数 $R(p)=\max p\cdot y\ \text{s.t.}\ F(y)=0$ & 都是给定价格下的最优值 \\[4pt]
$v$ 对 $(p,w)$ \textbf{零次齐次} & $R$ 对 $p$ \textbf{一次齐次} & 差别源于预算约束与技术约束的不同位置 \\[4pt]
$v$ 对 $p$ 非增 & $R$ 对 $p$ 非减 & 价格对消费者是成本，对生产者是收入 \\[4pt]
Roy 恒等式 $x_i=-\dfrac{\partial v/\partial p_i}{\partial v/\partial w}$ &
Hotelling 引理 $y_i=\dfrac{\partial R}{\partial p_i}$ & 都是包络定理的应用 \\[4pt]
支出函数 $e(p,\bar u)$ & 成本函数 $c(w,\bar y)$ & 互为另一侧的"最小支出/最小成本"值 \\[4pt]
希克斯需求 $h(p,\bar u)$ & 条件要素需求 $x(w,\bar y)$ & Shephard 引理：$\partial e/\partial p_i=h_i$，$\partial c/\partial w_j=x_j$ \\
\bottomrule
\end{tabular}
\end{table}

\subsection{与距离函数的接口（选读）}

Shephard 还给出过一种不依赖 $F=0$ 写法的技术描述：\textbf{产出距离函数}
\begin{equation}
D_O(y)=\min\bigl\{\theta>0:\ y/\theta\in Y\bigr\},
\label{eq:distance}
\end{equation}
其中 $Y=\{y:F(y)\le 0\}$ 为生产可能集。它满足：
$y\in Y$ 当且仅当 $D_O(y)\le 1$；$y$ 在前沿上当且仅当 $D_O(y)=1$；
$D_O$ 对 $y$ 一次齐次、非减、且在前沿附近与 $F$ 有相同的法向量方向
（相差一个正的比例因子）。因此，本案例中用 $\nabla F$ 给出的全部几何结论
（法向量、相切、$p=\lambda\nabla F$）都可以用 $\nabla D_O$ 原样复述一遍，
只是 $\lambda$ 的数值会随规范化方式改变，价格比等式
$MRT_{12}=p_1/p_2$ 则\textbf{完全不受影响}——它是不依赖任何规范化的经济内容。

\section{课堂讲授精简版}
\label{sec:lecture}

下面是一份可以直接照讲的口述稿（约 $2$ 分钟），随后是板书结构与两个提问。

\begin{quote}
\small
\textsl{我们举一个两产品的例子来理解收益最大化。
面包房一天的技术由转换函数 $F=y_1+\frac14 y_2^{2}-10=0$ 刻画：
等号成立时它在生产前沿上，小于零说明要素没用够，大于零则是做不到的组合。
市场价格是面包 $2$ 元、蛋糕 $4$ 元。
写下拉格朗日 $2y_1+4y_2-\lambda F$，一阶条件给出 $2=\lambda\cdot 1$、
$4=\lambda\cdot \frac{y_2}{2}$，于是 $\lambda=2$、$y_2=4$，回代前沿得 $y_1=6$，
最大收益 $28$ 元。
看它的经济含义：在最优点，价格比 $\frac{p_1}{p_2}=\frac12$ 恰好等于边际转换率
$F_1/F_2=\frac12$。$MRT_{12}$ 说的是"多做一个面包要放弃半个蛋糕"，
$p_1/p_2$ 说的是"市场上一个面包能换半个蛋糕"，
技术上的代价与市场上的所得正好打平，就没有再调整的余地了。
如果还没打平呢？请记住 $MRT_{12}$ 是成本、$p_1/p_2$ 是所得：
成本高于所得就要收缩这项产品，成本低于所得才去扩张。
几何上，$p=\lambda\nabla F$ 告诉我们价格向量与前沿的法向量共线，
$\lambda=2$ 是它们的比例；$-\lambda F$ 不是惩罚项，$\lambda$ 是影子值——
技术水平放松一单位，收益增加 $2$ 元，这也正是收益函数 $R=2c+8$ 的斜率。}
\end{quote}

\noindent\textbf{板书三行。}
\begin{enumerate}
  \item 约定：$F=0$（前沿）、$F<0$（可行）、$F>0$（不可行）；拉格朗日 $\mathcal{L}=p\cdot y-\lambda F$。
  \item 一阶条件与答案：$p_i=\lambda F_i\Rightarrow p_1/p_2=F_1/F_2=MRT_{12}$；
        $(y_1^{*},y_2^{*})=(6,4)$，$\lambda=2$，$\pi^{*}=28$。
  \item 三条含义：$MRT_{12}=\frac12$（技术代价$=$市场所得）、$p=2\nabla F$（法向量共线）、
        $\lambda=\partial R/\partial c=2$（影子值）。
\end{enumerate}

\noindent\textbf{两个课堂提问。}
\begin{enumerate}
  \item 保持技术不变，把两个价格同时翻倍，最优产出组合与最大收益如何变化？
        （答案：组合不变，仍为 $(6,4)$；收益翻倍为 $56$，因为 $R$ 对 $p$ 一次齐次。）
  \item 若当前产量为 $y_2=2$（前沿上的点），此时应当往哪个方向调整？为什么？
        （答案：$MRT_{12}=2/y_2=1>\frac12=p_1/p_2$，多做一个面包要放弃一个蛋糕、
        值 $4$ 元，而多卖一个面包只多得 $2$ 元，故应减少面包、增加蛋糕，
        即向 $y_2=4$ 的方向移动。）
\end{enumerate}

\section{本案例小结}
\label{sec:summary}

本案例从一条二阶多项式前沿出发，把生产者理论的收益侧完整走了一遍。

\begin{enumerate}
  \item \textbf{符号约定先行}。$F=0$ 为前沿、$F<0$ 为可行、$F>0$ 为不可行
        （\S\ref{sec:convention}）；约定必须与 $\partial F/\partial y_i>0$ 配对使用，
        中途不得更换，否则乘子符号与所有解释都会翻车。
  \item \textbf{一阶条件}。$\mathcal{L}=p\cdot y-\lambda F$，$p_i=\lambda F_i$，
        核心等式 $p_1/p_2=F_1/F_2=MRT_{12}$（\S\ref{sec:foc}）。
  \item \textbf{数值答案}。$(y_1^{*},y_2^{*})=(6,4)$，$\lambda=2$，$\pi^{*}=28$；
        消元法与沿前沿的收益对照给出两份交叉验证（\S\ref{sec:solve}）。
  \item \textbf{调整方向}。$MRT_{12}$ 是成本、$p_1/p_2$ 是所得：
        $MRT_{12}>p_1/p_2$ 时应\textbf{减少} $y_1$（\S\ref{sec:mrt}，
        附三份独立验证）。
  \item \textbf{几何}。$\nabla F$ 是前沿法向量，$p=\lambda\nabla F$ 表明二者共线同向，
        $\lambda$ 完成从"技术方向"到"市场价格"的缩放（\S\ref{sec:grad}）；
        切线斜率 $dy_1/dy_2=-F_2/F_1=-2$，与 $MRT$ 互为倒数或相反数，
        五种写法见表 \ref{tab:geometry}。
  \item \textbf{乘子的含义}。把技术水平写成参数 $c$ 后，$R=p_1c+p_2^{2}/p_1$，
        $\lambda=\partial R/\partial c=p_1=2$，既是影子值也是收益曲线的斜率
        （\S\ref{sec:shadow}）；$\partial R/\partial p_i=y_i^{*}$ 是 Hotelling 引理。
  \item \textbf{一个纠正}。$-\lambda F$ 不是惩罚项：$F<0$ 时它取值为正，
        真正的惩罚是二次对称的 $-\frac{\rho}{2}F^{2}$；
        把解推向前沿的机制是互补松弛条件（\S\ref{sec:penalty}）。
  \item \textbf{角点}。仅当 $p_2/p_1\le\sqrt{c}$ 时为内点，
        否则完全专业化：$y_2^{*}=\min\{2p_2/p_1,\,2\sqrt{c}\}$（\S\ref{sec:corner}）。
  \item \textbf{位置}。本案例的所有结论都能翻译到对偶体系中的标准位置：
        转换函数对偶于效用函数，收益函数对偶于间接效用函数，
        Hotelling 引理对偶于 Roy 恒等式（\S\ref{sec:duality}）。
\end{enumerate}

\subsection{检查清单}

解答同类题目时，逐项核对以下十条：

\begin{enumerate}
  \item 转换函数的三项约定写清楚了吗？（尤其可行的一侧）
  \item 拉格朗日中的符号与约定配套吗？（$-\lambda F$ 对应 $F\le 0$ 可行）
  \item 求出 $F_1,F_2$ 了吗？它们在最优点是否为正？
  \item 一阶条件写全了吗？（两个 $(y_i)$ 加一个乘子条件，共三个方程）
  \item $\lambda$ 是否为正？它与哪个价格、哪个梯度分量对应？
  \item 是否有不止一个解？（本例目标严格凹，唯一）
  \item 是否落在第一象限、是否在前沿上？（代回 $F=0$ 检验）
  \item $MRT$ 的方向核对了吗？（谁是分子、谁是分母、谁是负号）
  \item 是否需要与端点比较？（非凸技术、线性技术、极端价格）
  \item 经济解释齐了吗？（技术代价 $=$ 市场所得；法向量共线；乘子是影子值）
\end{enumerate}

\subsection{三个最常见的错误}

\begin{enumerate}
  \item \textbf{约定中途更换}。前半段用 $F\le 0$ 定义可行集，
        后半段却按 $-F\le 0$ 去解释，导致 $\lambda$ 变号、
        $MRT$ 差一个负号。防御办法：开篇写死约定，全篇不再讨论它的合理性。
  \item \textbf{$MRT$ 判据用反}。见 \S\ref{sec:mrt} 的注：
        $MRT_{12}>p_1/p_2$ 的正确应对是收缩 $y_1$，不是扩张它。
        判据的正确形式永远是 $d\pi=p_1\,dy_1+p_2\,dy_2$，
        把 $MRT$ 代入后再看符号，而不是凭"值钱"的直觉。
  \item \textbf{把 $-\lambda F$ 当惩罚}。等式约束问题没有惩罚环节；
        若需要"惩罚"，请明确 $\rho\to\infty$ 的罚函数，并注意它惩罚的是 $F^{2}$，
        且是对称的。
\end{enumerate}


\end{document}
